% !PRINTED 3
% Source chapter opening suppresses the running head and folio. Logical printed page 3 is
% fixed by the following explicit folio 4 and the preceding prelim sequence. The "1*" at
% the foot is a printer's signature, not a folio or body text, and is intentionally omitted.
\weberpage{3}

\begin{center}
\gesperrt{Erster Abschnitt.}
\end{center}

\begin{center}
\textbf{Die elliptischen Integrale.}
\end{center}

\begin{center}
\rule{3cm}{0.4pt}
\end{center}

\begin{center}
\textbf{\S\ 1. \gesperrt{Definition der elliptischen Integrale.}}
\end{center}

Wenn die systematische Darstellung der Integralrechnung bis zu dem Punkte gelangt ist, wo algebraische Integrale mit der Quadratwurzel aus einer Funktion ersten oder zweiten Grades auf Integrale rationaler Funktionen zurückgeführt werden, so tritt an dieser Stelle dem Lernenden eine Schranke entgegen, die er mit den ihm bis dahin zu Gebote stehenden Hilfsmitteln nicht zu übersteigen imstande ist. Das Streben nach einer Erweiterung der Hilfsmittel, um auch noch die nächste Klasse von Integralen der Forschung zugänglich zu machen, ist, wie es historisch der Anlaß gewesen, zu einem eingehenderen Studium der elliptischen Integrale und zur Einführung der elliptischen Funktionen, auch der naturgemäßeste und verständlichste Ausgangspunkt für den, der in die Theorie dieser Funktionen zuerst eingeführt werden soll. Es soll daher auch unsere nächste Aufgabe sein, uns mit den elliptischen Integralen und ihren wichtigsten Eigenschaften bekannt zu machen.

Die Definition eines elliptischen Integrals, von der wir ausgehen wollen, ist die folgende: Es bedeute $f(x)$ eine ganze rationale Funktion dritten oder vierten Grades mit vier verschiedenen Wurzeln
\weq{1}{f(x) = a_0 x^4 + 4 a_1 x^3 + 6 a_2 x^2 + 4 a_3 x + a_4,}
worin $a_0$ und $a_1$ nicht beide zugleich verschwinden; es sei ferner $\Phi(x, y)$ eine beliebige ganze oder gebrochene rationale Funktion der beiden Argumente $x, y$. Dann ist
\[
\int \Phi(x, \sqrt{f(x)})\, dx
\]
das allgemeine \gesperrt{elliptische Integral} und
\[
\Phi(x, \sqrt{f(x)})\, dx
\]
%CONT>>
